\documentclass[11pt,letterpaper]{article}

\usepackage{system/latex/curriculum}
\newcommand{\coursename}{Music Technology Engineering}
\hypersetup{pdftitle={Music Technology Engineering Mastery}}

\begin{document}

\coursecover
  {Music Technology Engineering Mastery}
  {Course III - An asynchronous, project-driven engineering curriculum}
  {Build working musical products that musicians enjoy using.}
  {Three pillars}
  {Music software | Physical instruments | Audio systems and infrastructure}
  {Learn. Retrieve. Predict. Build. Measure and listen. Review. Revisit.}
  {Trading Engineering and Oxide Systems Engineering are not prerequisites.}

\tableofcontents
\clearpage

\section{Mission}\label{mission}

Become an exceptionally compelling engineer for companies that build musical instruments, synthesizers, DJ equipment, digital audio workstations (DAWs), audio plugins, professional audio systems, and music software.

This curriculum covers \textbf{all nine engineering disciplines} identified in the career survey:

\begin{enumerate}
\item Embedded firmware for music equipment
\item Audio infrastructure and drivers
\item DAW and DJ application engineering
\item Digital signal processing (DSP)
\item Audio plugin development
\item Electronics and musical hardware engineering
\item Music software backends and platforms
\item Music AI and audio analysis
\item Professional audio and live-sound systems
\end{enumerate}

The aim is not merely to learn the technology, but to \textbf{build working musical products} that musicians enjoy using.

This course builds on existing experience in Go, C, backend systems, and embedded systems. A completed low-latency trading or Oxide systems curriculum would provide additional strength in C++, performance, concurrency, operating systems, and networking, but \textbf{neither is a prerequisite}. Reuse already-mastered material rather than repeating it.

\section{1. Philosophy and learning model}\label{philosophy-and-learning-model}

The target is a broadly capable music-technology engineer with \textbf{demonstrable depth}, not someone who has superficially sampled every discipline.

\subsection*{Three pillars}\label{three-pillars}

\begin{center}
\small
\begin{tabular}{>{\raggedright\arraybackslash}p{\dimexpr0.2800\linewidth-2\tabcolsep\relax}>{\raggedright\arraybackslash}p{\dimexpr0.7200\linewidth-2\tabcolsep\relax}}
\textbf{Pillar} & \textbf{Scope} \\
\textbf{A. Music software} & C++, real-time audio engines, DSP, synthesis, sampling, effects, plugins, DAWs, DJ software, desktop interfaces, music AI \\
\textbf{B. Physical instruments} & Embedded firmware, analog/digital electronics, MIDI, USB, controls, audio interfaces, sensors, PCBs, hardware/software integration \\
\textbf{C. Audio systems and infrastructure} & Drivers, operating-system audio, audio networking, synchronization, cloud services, collaboration, reliability, and product delivery \\
\end{tabular}
\end{center}

\subsection*{Required learning cycle}\label{required-learning-cycle}

Each major topic should follow this cycle:

\begin{enumerate}
\item \textbf{Learn:} Read or watch one focused explanation.
\item \textbf{Retrieve:} Close the material and explain the idea from memory.
\item \textbf{Predict:} State what an experiment should demonstrate.
\item \textbf{Build:} Implement the smallest meaningful example.
\item \textbf{Measure and listen:} Test the result technically and musically.
\item \textbf{Review:} Explain failures, surprises, and corrections.
\item \textbf{Revisit:} Repeat or extend the work later without relying on notes.
\end{enumerate}

Use retrieval practice, spaced repetition, interleaving, and repeated reconstruction. Avoid passive video/course consumption without implementation.

\subsection*{Progression rules}\label{progression-rules}

\begin{itemize}
\item \textbf{No fixed deadlines.} Each stage ends at an explicit mastery gate.
\item \textbf{Do not wait for perfect theory.} Start producing sounds and building usable tools early.
\item \textbf{Do not equate scope with excellence.} A small, excellent synth is better than an unfinished DAW clone.
\item \textbf{Don't postpone job applications} until all nine disciplines are mastered.
\item \textbf{Keep playing and making music.} Musical usefulness and tactile ergonomics are engineering requirements.
\end{itemize}

\section{2. Curriculum roadmap}\label{curriculum-roadmap}

\begin{center}
\small
\begin{tabular}{>{\raggedright\arraybackslash}p{\dimexpr0.0700\linewidth-2\tabcolsep\relax}>{\raggedright\arraybackslash}p{\dimexpr0.4300\linewidth-2\tabcolsep\relax}>{\raggedright\arraybackslash}p{\dimexpr0.5000\linewidth-2\tabcolsep\relax}}
\textbf{Stage} & \textbf{Topic} & \textbf{Primary deliverable} \\
0 & Music technology and product anatomy & Architecture breakdowns of three products \\
1 & Sound, music, mathematics, and digital audio & Audio signal and analysis experiments \\
2 & Real-time audio programming & \texttt{audio-engine} C++ playback/mixing engine \\
3 & DSP, synthesis, sampling, and effects & \texttt{synth-core} reusable DSP library \\
4 & Audio plugins and instrument applications & Polished VST3/AU software synthesizer \\
5 & DAW and DJ software engineering & \texttt{decklab} two-deck DJ application \\
6 & MIDI, controllers, and communication & Custom performance MIDI controller \\
7 & Embedded musical instruments & Standalone hardware synth/groovebox \\
8 & Audio electronics and hardware design & Custom, tested audio/control PCB \\
9 & Audio drivers and professional audio & \texttt{audio-transport-lab} and latency investigation \\
10 & Music platforms and cloud engineering & \texttt{sound-library} backend and synchronization \\
11 & Machine learning for musical applications & \texttt{music-assistant} evaluated intelligent features \\
12 & Integrated flagship product and open source & Finished public products and contributions \\
13 & Music-tech job preparation & Portfolio, demonstrations, and role-specific interview readiness \\
\end{tabular}
\end{center}

\textbf{Important:} This order is a learning dependency map, not a rule that you must finish everything before applying. Audio fundamentals and real-time programming come early; interview practice, documentation, user testing, and open source run alongside the main program.

\clearpage

\section{Stage 0 - Music technology and product anatomy}\label{stage-0-music-technology-and-product-anatomy}

\subsection*{Objective}\label{objective}

Understand how real music products are assembled from software, firmware, electronics, controls, audio processing, and user experience.

\subsection*{Study}\label{study}

Dissect the architecture of representative products:

\begin{itemize}
\item An Ableton-style DAW
\item A Pioneer DJ-style deck
\item A standalone groovebox
\item A digital hardware synthesizer
\item A MIDI controller
\item An audio interface
\item A software synth/plugin
\end{itemize}

For each, trace the \textbf{audio signal path}, the \textbf{control/data path}, and the \textbf{storage/preset path}. Identify which subsystems are application code, embedded firmware, DSP, device drivers, or electrical hardware.

\subsection*{Assignment}\label{assignment}

Produce \textbf{three product architecture breakdowns}, each containing:

\begin{itemize}
\item A subsystem diagram
\item The path from physical/user input to audible output
\item A description of hardware and software boundaries
\item Real-time requirements and likely failure modes
\item What you would prototype first if building a smaller version
\end{itemize}

\subsection*{Mastery gate}\label{mastery-gate}

You can explain the major subsystems of each product and why its design demands particular combinations of firmware, DSP, application software, and electronics.

\clearpage

\section{Stage 1 - Sound, music, mathematics, and digital audio}\label{stage-1-sound-music-mathematics-and-digital-audio}

\subsection*{Objective}\label{objective-1}

Develop an engineering understanding of how sound is represented, analyzed, generated, and manipulated.

\subsection*{Core topics}\label{core-topics}

\textbf{Physical and perceptual sound}

\begin{itemize}
\item Waveforms, frequency, amplitude, phase, and timbre
\item Harmonics and overtones
\item Loudness and decibels
\item Hearing and perceptual tradeoffs
\end{itemize}

\textbf{Digital audio}

\begin{itemize}
\item Sampling and the Nyquist principle
\item Quantization, bit depth, and sample formats
\item Aliasing and anti-aliasing
\item PCM, encoding, and digital audio files
\item Buffers and time-domain representation
\end{itemize}

\textbf{Mathematics for DSP}

\begin{itemize}
\item Sinusoids and complex numbers
\item Discrete-time signals
\item Convolution
\item Fourier transform and FFT
\item Frequency-domain representations
\item Basic difference equations and filter behavior
\end{itemize}

\textbf{Music fundamentals relevant to engineering}

\begin{itemize}
\item Intervals, scales, and chords
\item Tempo, meter, and rhythm
\item Envelopes, oscillators, filters, and modulation
\item Synthesis and sampling vocabulary
\end{itemize}

\subsection*{Experiments}\label{experiments}

\begin{enumerate}
\item Generate sine, square, sawtooth, and noise signals with code.
\item Build a waveform and spectrum viewer.
\item Demonstrate aliasing and quantization and listen to the artifacts.
\item Implement a simple low-pass filter.
\item Compare signals with the same fundamental frequency but different harmonic content.
\end{enumerate}

\subsection*{Mastery gate}\label{mastery-gate-1}

Given a common audible phenomenon, explain it in both time and frequency domains, recreate it in code, and test whether your explanation predicts the result.

\textbf{Anti-rabbit-hole boundary:} Learn derivations as needed to build and reason correctly. Do not defer implementation until you have completed an entire mathematics sequence.

\clearpage

\section{Stage 2 - Real-time audio programming}\label{stage-2-real-time-audio-programming}

\subsection*{Objective}\label{objective-2}

Write reliable C++ software that meets audio callback deadlines under real operating-system load.

\subsection*{Core topics}\label{core-topics-1}

\begin{itemize}
\item Audio blocks, buffers, sample rates, and channel layouts
\item Callback execution and deadlines
\item Audio thread versus UI/background threads
\item Buffer underflow/overflow, clicks, dropouts, and jitter
\item Priority inversion and thread coordination
\item Allocation avoidance and bounded work on the critical path
\item Lock-free or wait-free communication where justified
\item Parameter smoothing and denormal behavior
\item Core Audio and the architecture of ALSA, JACK/PipeWire, WASAPI, and ASIO
\item Instrumentation for missed audio deadlines
\end{itemize}

\subsection*{Project - \texttt{audio-engine}}\label{project-audio-engine}

Build a C++ engine that can:

\begin{itemize}
\item Play and mix samples
\item Route audio through processing nodes
\item Change parameters while audio is running
\item Handle start/stop and lifecycle changes
\item Expose audio diagnostics
\item Run under controlled artificial CPU/disk/UI load
\end{itemize}

Test callback behavior and missed deadlines, rather than merely measuring average throughput.

\subsection*{Mastery gate}\label{mastery-gate-2}

Your engine plays reliably while unrelated system activity occurs. You can explain how its threading and buffering prevent glitches, and you can measure the remaining limits.

\clearpage

\section{Stage 3 - DSP, synthesis, sampling, and effects}\label{stage-3-dsp-synthesis-sampling-and-effects}

\subsection*{Objective}\label{objective-3}

Design and implement the core sound-generation and sound-processing algorithms used in musical instruments.

\subsection*{Topic map}\label{topic-map}

\begin{center}
\small
\begin{tabular}{>{\raggedright\arraybackslash}p{\dimexpr0.3000\linewidth-2\tabcolsep\relax}>{\raggedright\arraybackslash}p{\dimexpr0.3400\linewidth-2\tabcolsep\relax}>{\raggedright\arraybackslash}p{\dimexpr0.3600\linewidth-2\tabcolsep\relax}}
\textbf{Synthesis} & \textbf{Effects} & \textbf{Analysis} \\
Subtractive & EQ and filters & FFT and spectrograms \\
Wavetable & Delay and reverb & Envelope detection \\
FM & Chorus and flanger & Pitch detection \\
Additive & Distortion and saturation & Transient/onset detection \\
Sampling & Compression and limiting & Tempo and rhythm analysis \\
\end{tabular}
\end{center}

Additionally study:

\begin{itemize}
\item Polyphony and voice allocation
\item Modulation matrices and parameter routing
\item Envelope generation and interpolation
\item Filter stability and frequency response
\item Oversampling and antialiasing
\item Numerical stability
\item SIMD and real-time DSP optimization when justified
\end{itemize}

\subsection*{Project - \texttt{synth-core}}\label{project-synth-core}

Create a reusable C++ synthesis/DSP library containing:

\begin{itemize}
\item Oscillators and noise sources
\item ADSR envelopes
\item Filters
\item Modulation
\item Polyphonic voice management
\item Basic sample playback
\item Delay and at least one other effect
\item Numerical tests and audio examples
\end{itemize}

Listen critically to the results. A numerically valid algorithm is not necessarily a musically useful one.

\subsection*{Mastery gate}\label{mastery-gate-3}

You can implement, test, profile, and explain common synthesis/effects algorithms. You can identify audible artifacts and propose defensible improvements.

\textbf{Optional tool:} \href{https://faust.grame.fr/}{Faust} for experimenting with DSP descriptions and generating implementations.

\clearpage

\section{Stage 4 - Audio plugins and instrument applications}\label{stage-4-audio-plugins-and-instrument-applications}

\subsection*{Objective}\label{objective-4}

Transform DSP code into a musician-facing product that works inside major DAWs.

\subsection*{Study}\label{study-1}

\begin{itemize}
\item Plugin lifecycles and host integration
\item VST3 and Audio Units; CLAP later
\item Host-to-plugin parameters and automation
\item MIDI event handling
\item Preset and project state persistence
\item Real-time-safe communication between UI and audio processing
\item C++ UI frameworks, especially JUCE
\item Cross-platform builds, packaging, versioning, and regression testing
\item Plugin validation and host compatibility
\end{itemize}

\subsection*{Flagship project A - Professional-quality software synthesizer}\label{flagship-project-a-professional-quality-software-synthesizer}

Build a complete, playable VST3/AU instrument using \texttt{synth-core}.

Required features:

\begin{itemize}
\item Oscillators, filters, envelopes, and modulation
\item Polyphony and velocity response
\item A coherent, responsive graphical UI
\item Preset management
\item Parameter automation
\item Reliable save/reopen behavior in a DAW
\item A stable, glitch-resistant audio path
\item Installation instructions and a short product demonstration
\end{itemize}

Validate in real hosts; use tools such as \href{https://github.com/Tracktion/pluginval}{pluginval}.

\subsection*{Mastery gate}\label{mastery-gate-4}

A musician can install the plugin, open it in Ableton, produce a useful sound, automate it, save the project, reopen it, and hear the same result reliably.

\textbf{Application milestone:} This is the first major readiness point for audio-plugin, software-instrument, and some audio-engine positions.

\textbf{Primary resources:} \href{https://juce.com/learn/course/}{JUCE course} and \href{https://steinbergmedia.github.io/vst3_dev_portal/}{VST3 Developer Portal}.

\clearpage

\section{Stage 5 - DAW and DJ software engineering}\label{stage-5-daw-and-dj-software-engineering}

\subsection*{Objective}\label{objective-5}

Understand the architecture behind timeline-based creative tools and DJ applications.

\subsection*{Study}\label{study-2}

\textbf{DAW/app architecture}

\begin{itemize}
\item Audio routing graphs and buses
\item Multitrack mixing and sample-accurate scheduling
\item Timelines, sequencing, and playback state
\item Non-destructive editing, undo/redo, and persistence
\item Plugin hosting
\item Waveform visualization and desktop UI architecture
\end{itemize}

\textbf{DJ-specific systems}

\begin{itemize}
\item Beat grids, BPM analysis, cue points, and looping
\item Pitch shifting and time stretching
\item Tempo sync and beat alignment
\item Jog-wheel and physical-controller behavior
\item DJ library metadata and browsing
\item Real-time FX chains
\end{itemize}

\subsection*{Project - \texttt{decklab}}\label{project-decklab}

Build a compact two-deck DJ application with:

\begin{itemize}
\item Playback, cueing, looping, and mixing
\item Tempo and pitch controls
\item Waveform visualization
\item Track library/metadata
\item MIDI mapping
\item A basic effects chain
\item Basic beat synchronization
\end{itemize}

Test with physical DJ hardware when feasible. Study the architecture and contribution process of \href{https://github.com/mixxxdj/mixxx}{Mixxx}.

\subsection*{Mastery gate}\label{mastery-gate-5}

Perform a short real DJ mix with the application. Explain its playback architecture, synchronization behavior, and important limitations compared with commercial products.

\textbf{Scope control:} Do not attempt to clone every feature in Ableton Live or rekordbox.

\clearpage

\section{Stage 6 - MIDI, controllers, and musical communication}\label{stage-6-midi-controllers-and-musical-communication}

\subsection*{Objective}\label{objective-6}

Create expressive physical controls and reliable communication between musical devices and applications.

\subsection*{Study}\label{study-3}

\begin{itemize}
\item MIDI 1.0 notes, velocity, CC, pitch bend, clock, and SysEx
\item MIDI over USB
\item MIDI 2.0 and Universal MIDI Packets
\item Device discovery and mappings
\item Timing, synchronization, and jitter
\item OSC
\item BLE MIDI concepts
\item Ableton Link and network synchronization
\item Encoders, buttons, potentiometers, pads, and LED feedback
\item Human factors and physical interaction design
\end{itemize}

\subsection*{Project - Custom performance MIDI controller}\label{project-custom-performance-midi-controller}

Build a hardware controller with physical controls, firmware, persistent mapping, and feedback.

Integrate it with your software synthesizer and DJ application. Consider a small display, LED rings, and remappable controls if they improve real musical use.

\subsection*{Mastery gate}\label{mastery-gate-6}

The device works predictably during performance, persists mappings, communicates with software reliably, and can be used without detailed firmware knowledge.

\textbf{Reference:} \href{https://midi.org/midi-2-0}{MIDI Association} and \href{https://www.ableton.com/en/link/}{Ableton Link}.

\clearpage

\section{Stage 7 - Embedded musical instruments}\label{stage-7-embedded-musical-instruments}

\subsection*{Objective}\label{objective-7}

Create a standalone digital instrument that makes sound without a host computer.

\subsection*{Study}\label{study-4}

\begin{itemize}
\item Embedded C/C++ audio architectures
\item Microcontroller, DSP, and SoC tradeoffs
\item Audio codecs, ADC/DAC, and audio output
\item I\textsuperscript{2}S/SAI interfaces
\item DMA and interrupt-driven sample processing
\item Fixed-size audio buffers and real-time scheduling
\item MCU RAM/flash limitations
\item Preset storage, boot behavior, watchdogs, and firmware updates
\item Portable DSP versus platform-specific optimizations
\end{itemize}

Use a practical audio development platform, such as \href{https://docs.daisy.audio/}{Electrosmith Daisy} or \href{https://learn.bela.io/tutorials/}{Bela}.

\subsection*{Flagship project B - Standalone hardware instrument}\label{flagship-project-b-standalone-hardware-instrument}

Port suitable pieces of \texttt{synth-core} and implement:

\begin{itemize}
\item Audio generation or sample playback
\item Physical inputs and an understandable UI
\item Preset storage
\item MIDI support
\item Basic sequencing
\item Reliable startup, recovery, and firmware updates
\item Measurement of audio latency and resource use
\end{itemize}

The hardware can be a compact synth, sampler, or groovebox. Its precise form should follow what you most enjoy making music with.

\subsection*{Mastery gate}\label{mastery-gate-7}

The instrument boots reliably, operates over extended sessions, responds expressively, preserves settings, and has a documented timing/resource profile.

\textbf{Application milestone:} With Stage 6 and some hands-on hardware debugging, this is a strong foundation for embedded musical-hardware firmware applications.

\clearpage

\section{Stage 8 - Audio electronics and hardware design}\label{stage-8-audio-electronics-and-hardware-design}

\subsection*{Objective}\label{objective-8}

Become capable of designing and debugging the analog/digital electronics surrounding audio software.

\subsection*{Study}\label{study-5}

\begin{itemize}
\item Voltage, current, impedance, and signal levels
\item Op-amps and analog filter circuits
\item Grounding, shielding, and noise
\item ADC/DAC fundamentals
\item Clocking and converter behavior
\item Power regulation and analog/digital isolation concepts
\item Connectors and protection
\item PCB schematics, layout, manufacturing files, and bring-up
\item Measurement with multimeters, oscilloscopes, and audio test signals
\item Electromagnetic compatibility at a practical level
\end{itemize}

\subsection*{Project - Custom audio/control PCB}\label{project-custom-audiocontrol-pcb}

Use \href{https://www.kicad.org/}{KiCad} to design a low-voltage PCB for your instrument.

Start with a simple control/interface board, then increase complexity only when justified. Deliver:

\begin{itemize}
\item Schematic and PCB layout
\item Bill of materials
\item Assembly and bring-up instructions
\item Test plan and test results
\item Measured electrical behavior
\item Diagnosis of at least one real fault or design revision
\end{itemize}

\subsection*{Mastery gate}\label{mastery-gate-8}

You can read component datasheets, trace analog/digital signal paths, debug board-level issues, and bring up a manufactured PCB.

\textbf{Safety:} Begin with low-voltage, USB-powered systems. Avoid mains-powered circuit design without appropriate expertise and safety controls.

\textbf{Career qualification:} This gives meaningful hardware/product fluency; it does not substitute for deep EE experience in specialist analog/RF hardware design.

\clearpage

\section{Stage 9 - Audio interfaces, drivers, and professional audio}\label{stage-9-audio-interfaces-drivers-and-professional-audio}

\subsection*{Objective}\label{objective-9}

Understand audio from software callback to operating system, interface hardware, and professional audio network.

\subsection*{Study}\label{study-6}

\begin{itemize}
\item OS audio I/O stacks and device enumeration
\item USB Audio Class concepts
\item Sample-rate conversion and clock-domain crossings
\item Buffering, latency compensation, and device synchronization
\item Drivers and firmware/host interfaces
\item Multichannel routing and digital mixers
\item Networked audio and audio-over-IP
\item Packet timing, jitter buffers, dropped packets, and recovery
\item AES67 concepts and professional network synchronization
\item Redundancy, monitoring, and failure diagnosis
\end{itemize}

\subsection*{Project - \texttt{audio-transport-lab}}\label{project-audio-transport-lab}

Stream low-latency multichannel audio between two computers.

Implement and analyze:

\begin{itemize}
\item Capture and playback
\item Packetization and transport
\item Jitter buffering and clock drift
\item Packet loss/disconnection recovery
\item Measurements of end-to-end latency and audio quality
\item Diagnostic tools/logs for identifying glitches
\end{itemize}

\subsection*{Mastery gate}\label{mastery-gate-9}

You can isolate and explain timing failures across application, driver, network, and device boundaries, with supporting measurements.

\clearpage

\section{Stage 10 - Music platforms and cloud engineering}\label{stage-10-music-platforms-and-cloud-engineering}

\subsection*{Objective}\label{objective-10}

Apply existing backend and distributed-systems strengths to cloud capabilities used by music creators and connected devices.

\subsection*{Study}\label{study-7}

\begin{itemize}
\item Preset libraries and sample catalogs
\item Search and metadata models
\item Cloud storage and asset lifecycle
\item User accounts, permissions, licensing, and entitlements
\item Cross-device synchronization and conflict resolution
\item Collaboration and sharing
\item Device configuration and update services
\item Event-driven workflows and observability
\item Music-product UX implications for sync and offline operation
\end{itemize}

\subsection*{Project - \texttt{sound-library}}\label{project-sound-library}

Build a backend for:

\begin{itemize}
\item Preset/sample metadata
\item Search and user collections
\item Synchronization between plugin, controller, and instrument
\item Versioned presets and schema evolution
\item Sharing/collaboration where useful
\item Secure APIs and offline-recovery behavior
\end{itemize}

Use Go unless a different language provides a clear technical benefit.

\subsection*{Mastery gate}\label{mastery-gate-10}

Changes propagate predictably across applications/devices, conflicts are resolved deliberately, and tests demonstrate recovery under failure.

\textbf{Application milestone:} Existing backend experience may already qualify you for music-platform positions. Do not postpone these applications for the DSP/hardware curriculum.

\clearpage

\section{Stage 11 - Machine learning for musical applications}\label{stage-11-machine-learning-for-musical-applications}

\subsection*{Objective}\label{objective-11}

Implement, evaluate, and productize useful intelligent music features.

\subsection*{Study}\label{study-8}

\begin{itemize}
\item Onset, beat, tempo, and key estimation
\item Pitch and harmonic analysis
\item Timbre and audio embeddings
\item Musical similarity and retrieval
\item Source separation and transcription concepts
\item Conventional DSP versus ML approaches
\item Inference latency and resource limits
\item Model evaluation and dataset quality
\item Copyright, licensing, provenance, and responsible model use
\end{itemize}

Begin with classical DSP and audio analysis. Use Python for experimentation, and integrate mature libraries/models when appropriate rather than treating foundation-model training as a prerequisite.

\subsection*{Project - \texttt{music-assistant}}\label{project-music-assistant}

Implement \textbf{at least two} features, such as:

\begin{itemize}
\item BPM/key estimation for \texttt{decklab}
\item Similar-sound sample search
\item Section/transient detection
\item Stem separation using an existing model
\item Preset or instrument recommendations
\end{itemize}

Create reproducible evaluation data. Compare predicted versus measured behavior and include bad cases, latency, and limits.

\subsection*{Mastery gate}\label{mastery-gate-11}

You can defend the algorithm, evaluation methodology, resource costs, and practical musical value of each feature.

\textbf{Career qualification:} Deep audio-ML research would require an additional advanced specialization; the present gate targets solid audio-ML software engineering.

\textbf{Resources:} \href{https://librosa.org/doc/latest/auto_tutorials/}{librosa tutorials} and \href{https://essentia.upf.edu/}{Essentia}.

\clearpage

\section{Stage 12 - Integrated flagship products and open source}\label{stage-12-integrated-flagship-products-and-open-source}

\subsection*{Objective}\label{objective-12}

Turn the course's projects into an unusually convincing public body of music-technology engineering work.

\subsection*{Preferred integration architecture}\label{preferred-integration-architecture}

{\footnotesize
\begin{verbatim}
Physical MIDI controller
          |
          v
+-------------------------+
| Software synth / plugin |<---- shared DSP concepts ----> Hardware synth
+-------------------------+
          |
          v
     DJ app / mixer
          |
          v
Preset library / sync / metadata
          |
          v
Audio analysis and smart features
\end{verbatim}
}

Keep components modular. They do not need to become one huge monolithic product.

\subsection*{Four flagship portfolio artifacts}\label{four-flagship-portfolio-artifacts}

\textbf{1. A polished synthesizer plugin}

Demonstrates C++, DSP, real-time programming, plugins, desktop UX, and musical judgment.

\textbf{2. A standalone groovebox/synth}

Demonstrates embedded firmware, audio interfaces, MIDI, electrical design, and physical interaction.

\textbf{3. A DJ application with controller integration}

Demonstrates real-time desktop software, playback architecture, synchronization, and creative workflows.

\textbf{4. An audio-systems/music-platform laboratory}

Demonstrates backend systems, device sync, audio networking, reliability, and music analysis.

\subsection*{Quality expectations}\label{quality-expectations}

For each flagship artifact, publish:

\begin{itemize}
\item Working code and an understandable build/install procedure
\item Automated tests and realistic manual tests
\item Architecture diagrams and major design decisions
\item Measurements where relevant
\item A short video or audio demonstration
\item Limitations and known failures
\item User feedback and improvements
\end{itemize}

\clearpage
\subsection*{Open-source contribution track}\label{open-source-contribution-track}

Study and contribute to real music/audio codebases when useful:

\begin{itemize}
\item \href{https://github.com/mixxxdj/mixxx}{Mixxx}
\item \href{https://github.com/surge-synthesizer/surge}{Surge XT}
\item \href{https://github.com/Chowdhury-DSP}{ChowDSP}
\item \href{https://github.com/Tracktion/pluginval}{pluginval}
\item \href{https://github.com/MTG/essentia}{Essentia}
\end{itemize}

Start by building a project, running tests, understanding a subsystem, reproducing a bug, and submitting a tested change. Respond to review thoughtfully.

\textbf{Mastery gate:} You have multiple polished public artifacts, a technical narrative you can defend, and ideally at least one substantive externally reviewed contribution.

\clearpage

\section{Stage 13 - Music-technology interviewing and career preparation}\label{stage-13-music-technology-interviewing-and-career-preparation}

\subsection*{Objective}\label{objective-13}

Translate your technical and musical work into credible applications and realistic interview performance.

\subsection*{Interview preparation by role}\label{interview-preparation-by-role}

\begin{center}
\small
\begin{tabular}{>{\raggedright\arraybackslash}p{\dimexpr0.2700\linewidth-2\tabcolsep\relax}>{\raggedright\arraybackslash}p{\dimexpr0.7300\linewidth-2\tabcolsep\relax}}
\textbf{Role} & \textbf{Emphasize} \\
Audio DSP & Filters, DSP math, sampling, FFT, numerical stability, C++ \\
Plugin/DAW engineer & C++, audio threading, plugin APIs, GUI/audio synchronization, architecture \\
Embedded firmware & C/C++, interrupts, DMA, I\textsuperscript{2}S, device debugging, deterministic timing \\
Hardware engineer & Circuit analysis, test equipment, PCB design, audio signal integrity \\
Professional audio & Drivers, audio I/O, networking, clocks, fault tolerance \\
Backend/platform & Distributed systems, API/data design, reliability, scalability \\
Music AI & Signal analysis, evaluation, inference, ML tradeoffs \\
\end{tabular}
\end{center}

\subsection*{Practice formats}\label{practice-formats}

\begin{itemize}
\item Timed coding exercises
\item Short DSP implementation challenges
\item Recorded design explanations
\item Debugging sessions with deliberate faults
\item Product walkthroughs
\item Async code reviews from experienced engineers
\item Occasional external mock interviews, without making scheduling a learning dependency
\end{itemize}

\subsection*{Portfolio demonstration standard}\label{portfolio-demonstration-standard}

For every substantial project, prepare to show:

\begin{enumerate}
\item A concise demonstration of its musical use
\item The architecture
\item One difficult engineering problem
\item A bug or failure investigation
\item A measurable tradeoff or improvement
\item What you would change in a production version
\end{enumerate}

\subsection*{Mastery gate}\label{mastery-gate-12}

You can interview credibly for at least one chosen specialty and present compelling, technically defensible examples of your work.

\clearpage

\section{3. Application gates: do not wait for the whole course}\label{application-gates-do-not-wait-for-the-whole-course}

\begin{center}
\small
\begin{tabular}{>{\raggedright\arraybackslash}p{\dimexpr0.3600\linewidth-2\tabcolsep\relax}>{\raggedright\arraybackslash}p{\dimexpr0.6400\linewidth-2\tabcolsep\relax}}
\textbf{Target} & \textbf{Reasonable point to begin applying} \\
Music backend/platform & Potentially now, with targeted domain preparation \\
Audio engine/DSP/plugin & After Stages 1-4, with a polished plugin \\
DAW/DJ application engineer & After Stages 1-5, with \texttt{decklab} \\
MIDI/controller firmware & After Stages 6-7, supported by real-time audio basics \\
Instrument firmware & After Stage 7 and hands-on hardware debugging \\
Hardware engineer & After Stage 8 and a functioning PCB; deeper EE may still be required \\
Professional audio systems & After Stage 9, with networking and diagnostics evidence \\
Music AI engineer & After Stage 11 and a reproducible evaluated ML project \\
Cross-disciplinary product engineer & After multiple finished, integrated portfolio projects \\
\end{tabular}
\end{center}

These are \textbf{learning milestones}, not guarantees that a particular employer will consider you qualified.

\subsection*{Three completion levels}\label{three-completion-levels}

\textbf{Level 1 - Employable specialty candidate}

A relevant functioning project, knowledge of the role's core engineering, and credible interview readiness.

\textbf{Level 2 - Strong cross-disciplinary engineer}

Meaningful software, embedded, DSP, electronics, and audio-system experience with well-documented public work.

\textbf{Level 3 - Exceptionally compelling specialist with broad product fluency}

One or more musically polished products, unusual depth in selected disciplines, external reviews/contributions, and ability to defend architectural choices throughout the audio stack.

Level 3 is the ambitious target. It does \textbf{not} mean being simultaneously equivalent to a lifelong analog-circuit specialist, senior DSP researcher, embedded architect, and principal cloud engineer.

\section{4. Product testing is part of the curriculum}\label{product-testing-is-part-of-the-curriculum}

Music technology is not judged only by benchmarks or test suites. Continue DJing, producing music, and experimenting with instruments during the entire course.

For each major project:

\begin{itemize}
\item Use it in a real session, not only a scripted demo.
\item Ask another musician to try it without guidance.
\item Note confusing controls, poor ergonomics, sound-quality weaknesses, and workflow friction.
\item Treat artistic and usability feedback as inputs to the engineering backlog.
\end{itemize}

A perfectly stable plugin that musicians find frustrating is not an exceptional product. A controller with excellent firmware and poor tactile usability is similarly incomplete.

\clearpage
\section{5. Resource stack}\label{resource-stack}

Use a limited set of primary resources and add specialized materials only as they become necessary.

\begin{center}
\small
\begin{tabular}{>{\raggedright\arraybackslash}p{\dimexpr0.4800\linewidth-2\tabcolsep\relax}>{\raggedright\arraybackslash}p{\dimexpr0.5200\linewidth-2\tabcolsep\relax}}
\textbf{Subject} & \textbf{Starting resource} \\
Audio plugin/real-time C++ & \href{https://juce.com/learn/course/}{JUCE course} \\
DSP and synthesis experimentation & \href{https://faustdoc.grame.fr/}{Faust documentation} \\
Embedded audio & \href{https://learn.bela.io/tutorials/}{Bela tutorials} \\
Hardware synth firmware & \href{https://docs.daisy.audio/}{Daisy documentation} \\
MIDI & \href{https://midi.org/midi-2-0}{MIDI Association} \\
VST3 & \href{https://steinbergmedia.github.io/vst3_dev_portal/}{Steinberg Developer Portal} \\
PCB design & \href{https://www.kicad.org/help/learning-resources/}{KiCad resources} \\
Music and audio analysis & \href{https://librosa.org/doc/latest/auto_tutorials/}{librosa tutorials} \\
Real DJ application source & \href{https://github.com/mixxxdj/mixxx}{Mixxx} \\
Real synthesizer source & \href{https://github.com/surge-synthesizer/surge}{Surge XT} \\
Plugin validation & \href{https://github.com/Tracktion/pluginval}{pluginval} \\
Audio-information retrieval & \href{https://essentia.upf.edu/}{Essentia} \\
\end{tabular}
\end{center}

\section{6. Anti-rabbit-hole rules}\label{anti-rabbit-hole-rules}

\begin{enumerate}[leftmargin=*,widest=10]
\item \textbf{Don't clone a commercial product.} Build a narrowly scoped instrument or creative tool that works well.
\item \textbf{Don't separate theory from sound.} A math concept should lead quickly to audible implementation.
\item \textbf{Don't design custom boards too early.} Prototype on desktop and inexpensive development hardware first.
\item \textbf{Don't chase exotic DSP before the basics sound good.} Musical quality and robustness matter.
\item \textbf{Don't reimplement standards and drivers unnecessarily.} Use established interfaces where possible.
\item \textbf{Don't confuse technical success with product usefulness.} Incorporate musician feedback.
\item \textbf{Don't let breadth block applications.} Apply when a specialty-specific gate is satisfied.
\item \textbf{Don't collect repositories.} Polish the strongest products and make them easy to inspect.
\item \textbf{Don't optimize without a measurement or audible reason.} Preserve correctness and sound quality.
\item \textbf{Don't buy expensive hardware without a validated need.} Start with simulation and affordable development boards.
\end{enumerate}

\section{7. Relationship to the quant and Oxide courses}\label{relationship-to-the-quant-and-oxide-courses}

\begin{center}
\small
\begin{tabular}{>{\raggedright\arraybackslash}p{\dimexpr0.3100\linewidth-2\tabcolsep\relax}>{\raggedright\arraybackslash}p{\dimexpr0.2300\linewidth-2\tabcolsep\relax}>{\raggedright\arraybackslash}p{\dimexpr0.2300\linewidth-2\tabcolsep\relax}>{\raggedright\arraybackslash}p{\dimexpr0.2300\linewidth-2\tabcolsep\relax}}
\textbf{Capability} & \textbf{Quant trading} & \textbf{Oxide systems} & \textbf{Music technology} \\
C++ systems programming & Essential & Useful & Essential \\
Performance engineering & Essential & Essential & Essential \\
Concurrency / real time & Essential & Essential & Essential \\
Operating systems & Important & Essential & Important \\
Networking & Important & Essential & Important \\
Embedded systems & Occasional & Important & Essential for physical products \\
DSP & Limited & Specialized & Essential \\
Electronics & Limited & Important & Essential for hardware roles \\
Cloud/distributed systems & Important & Essential & Useful \\
Technical product design & Useful & Important & Essential \\
\end{tabular}
\end{center}

These curricula share a strong systems-programming foundation. This third course adds a new domain: \textbf{the physics, mathematics, perception, and interaction design of sound and musical instruments}.

\section{8. Recommended first move}\label{recommended-first-move}

\textbf{Build a small but polished software synthesizer plugin first.}

It is inexpensive to develop, immediately useful in Ableton, and introduces the key audio/DSP/real-time skills without buying or designing hardware. Later, porting the same synthesis concepts to a standalone embedded instrument creates a natural bridge to MIDI controllers, circuit design, and physical musical products.

The long-term goal is broad capability. The near-term goal is to \textbf{finish one excellent thing, learn from using it, and let that success finance the next layer of skill}.

\end{document}
